\subsection{$1/x$ 的主值}
\label{sec:ch3-principal-value-one-over-x}

我们定义一个称为 $1/x$ 的主值的缓增分布, 记作 $\operatorname{p.v.}\,1/x$:
\[
 \operatorname{p.v.}\frac1x(\varphi)
 =\lim_{\varepsilon\to0}\int_{|x|>\varepsilon}\frac{\varphi(x)}{x}\,dx,
 \qquad \varphi\in\mathcal S.
\]
为了看出这个表达式定义了一个缓增分布, 将它改写为
\[
 \operatorname{p.v.}\frac1x(\varphi)
 =\int_{|x|<1}\frac{\varphi(x)-\varphi(0)}{x}\,dx
  +\int_{|x|>1}\frac{\varphi(x)}{x}\,dx;
\]
这是因为 $1/x$ 在 $\varepsilon<|x|<1$ 上的积分为零. 于是立即得到
\[
 \left|\operatorname{p.v.}\frac1x(\varphi)\right|
 \le C\bigl(\|\varphi'\|_\infty+\|x\varphi\|_\infty\bigr).
\]

\begin{Proposition}
\label{prop:ch3-conjugate-poisson-principal-value}
在 $\mathcal S'$ 中,
\[
 \lim_{t\to0}Q_t=\frac1\pi\operatorname{p.v.}\frac1x.
\]
\end{Proposition}

\begin{Proof}
对每个 $\varepsilon>0$, 函数 $\psi_\varepsilon(x)=x^{-1}\chi_{\{|x|>\varepsilon\}}$ 有界并定义缓增分布. 由定义立即知道, 在 $\mathcal S'$ 中
\[
 \lim_{\varepsilon\to0}\psi_\varepsilon=\operatorname{p.v.}\frac1x.
\]
因此只需证明在 $\mathcal S'$ 中
\[
 \lim_{t\to0}\left(Q_t-\frac1\pi\psi_t\right)=0.
\]

\hypertarget{chapter-03-page-063}{}
固定 $\varphi\in\mathcal S$. 则
\[
\begin{aligned}
 (\pi Q_t-\psi_t)(\varphi)
 &=\int_{\mathbb R}\frac{x\varphi(x)}{t^2+x^2}\,dx
   -\int_{|x|>t}\frac{\varphi(x)}x\,dx\\
 &=\int_{|x|<t}\frac{x\varphi(x)}{t^2+x^2}\,dx
   +\int_{|x|>t}\left(\frac{x}{t^2+x^2}-\frac1x\right)\varphi(x)\,dx\\
 &=\int_{|x|<1}\frac{x\varphi(tx)}{1+x^2}\,dx
   -\int_{|x|>1}\frac{\varphi(tx)}{x(1+x^2)}\,dx.
\end{aligned}
\]
令 $t\to0$ 并应用控制收敛定理, 得到两个在对称区域上对奇函数的积分. 因而极限为零.
\end{Proof}

由命题~\ref{prop:ch3-conjugate-poisson-principal-value} 可得
\[
 \lim_{t\to0}Q_t*f(x)
 =\frac1\pi\lim_{\varepsilon\to0}\int_{|y|>\varepsilon}\frac{f(x-y)}y\,dy,
\]
再由 Fourier 变换在 $\mathcal S'$ 上的连续性及\eqref{eq:ch3-conjugate-poisson-fourier}, 得
\[
 \left(\frac1\pi\operatorname{p.v.}\frac1x\right)^{\wedge}(\xi)
 =-i\operatorname{sgn}(\xi).
\]

给定 $f\in\mathcal S$, 用下列任一等价表达式定义它的 Hilbert 变换:
\[
 Hf=\lim_{t\to0}Q_t*f,
 \qquad
 Hf=\frac1\pi\operatorname{p.v.}\frac1x*f,
 \qquad
 \widehat{Hf}(\xi)=-i\operatorname{sgn}(\xi)\widehat f(\xi).
\]
第三个表达式还允许我们定义 $L^2(\mathbb R)$ 中函数的 Hilbert 变换; 它满足
\begin{equation}
\label{eq:ch3-hilbert-l2-isometry}
 \|Hf\|_2=\|f\|_2,
\end{equation}
\begin{equation}
\label{eq:ch3-hilbert-square-minus-identity}
 H(Hf)=-f,
\end{equation}
以及
\begin{equation}
\label{eq:ch3-hilbert-skew-adjoint}
 \int Hf\cdot g=-\int f\cdot Hg.
\end{equation}
